\subsection{Factorisation à travers l’opérateur constitutif}

Conservons la coordinatisation angulaire relevée de la
section~\ref{iii:sec:hojas}. Pour rendre explicite la conversion angulaire,
écrivons
\begin{equation}
 A_{\deg}=\frac{180x}{\pi},\qquad
 C^*_{\deg}=\frac{180y}{\pi},\qquad
 q_\pm=e^{-x\mp y},\qquad s_\pm=q_\pm^{30}.
 \label{iii:eq:barbero-angular-lift}
\end{equation}
L’orientation considérée dans \eqref{iii:eq:rchannels} correspond à
$0<y<x$ ; l’échange de feuillets s’exprime par $y\mapsto-y$ sur le
domaine plus large $x>|y|$. Les exponentielles réelles agissent sur ces
relèvements, qui conservent les valeurs de $x\pm y$ pendant la
composition.

La fonctionnelle orientée de Barbero--Immirzi est
\begin{equation}
 \gamma=\frac{\pi A_{\deg}C^*_{\deg}}{180}
   +12\bigl[S_{90}(q_-)-S_{90}(q_+)\bigr]
   -\bigl[S_{120}(q_-)-S_{120}(q_+)\bigr].
 \label{iii:eq:barbero-functional}
\end{equation}
La première composante est l’évaluation angulaire ; la seconde est la
différence du lecteur de la chaîne fondamentale dans les canaux conjugués.
Les arguments sont reçus du transport angulaire précédemment construit.
La formule détermine un scalaire adimensionnel sur cette structure de
départ.

Soit $\psi:(3/4,1)\to(0,1)$ l’inverse de $f$ construit au moyen de
\eqref{iii:eq:cubic}. Définissons
\begin{equation}
 \mathcal H(s)=\frac{12s^3}{1-s^9}
             -\frac{s^4}{1-s^{12}}
             -\frac{(\log s)^2}{20\pi},\qquad 0<s<1.
 \label{iii:eq:barbero-h}
\end{equation}
Cette fonction est analytique réelle sur son domaine. Les exposants $3,4,9,12$
résultent de l’expression des hauteurs $90,120$ et des retours $270,360$ en
unités de $30$ ; le coefficient de la partie logarithmique est fixé lorsqu’on
transporte la contribution angulaire.

\begin{theorem}[Expression constitutive orientée de Barbero--Immirzi]
\label{iii:thm:barbero-factorization}
Dans la coordinatisation \eqref{iii:eq:barbero-angular-lift},
\begin{equation}
 \gamma=\mathcal H(s_-)-\mathcal H(s_+)
       =\mathcal H\bigl(\psi(r_-)\bigr)
        -\mathcal H\bigl(\psi(r_+)\bigr).
 \label{iii:eq:barbero-factorization}
\end{equation}
Sur la représentation bidimensionnelle des canaux,
\begin{equation}
 \gamma=-\operatorname{Tr}
    \left[\mathcal R\,\mathcal H\bigl(\psi(\mathcal V)\bigr)\right],
 \label{iii:eq:barbero-trace}
\end{equation}
où $\operatorname{Tr}$ est la trace ordinaire, avec
$\operatorname{Tr}I=2$.
\end{theorem}
\begin{proof}
La substitution $s=q^{30}$ dans \eqref{iii:eq:barbero-return-series} donne
\[
 S_{90}(q)=\frac{s^3}{1-s^9},\qquad
 S_{120}(q)=\frac{s^4}{1-s^{12}}.
\]
De plus, \eqref{iii:eq:barbero-angular-lift} implique
\[
 \log s_\pm=-30(x\pm y)
 =-\frac\pi6(A_{\deg}\pm C^*_{\deg}).
\]
Par différence de carrés,
\begin{align}
 \frac{(\log s_+)^2-(\log s_-)^2}{20\pi}
 &=\frac{900\bigl[(x+y)^2-(x-y)^2\bigr]}{20\pi}\\
 &=\frac{180xy}{\pi}
 =\frac{\pi A_{\deg}C^*_{\deg}}{180}.
 \label{iii:eq:barbero-log-term}
\end{align}
La somme de ces identités prouve la première égalité de
\eqref{iii:eq:barbero-factorization}. La seconde utilise l’inversion
$\psi(f(s_\pm))=s_\pm$ démontrée dans la
section~\ref{iii:sec:constitutiva}.

Pour l’expression opératorielle, les projecteurs de feuillet sont de rang un et
$\mathcal RP_\pm=\pm P_\pm$. Le calcul spectral fournit
\[
 \mathcal H\bigl(\psi(\mathcal V)\bigr)
 =\mathcal H(s_+)P_++\mathcal H(s_-)P_-.
\]
Multiplier par $\mathcal R$ et prendre la trace ordinaire donne
$\mathcal H(s_+)-\mathcal H(s_-)$. Son opposé est $\gamma$.
\end{proof}

La marque d’orientation $\mathcal R$ fait partie de la fonctionnelle. Si
l’on maintient cette marque fixe et que l’on échange les canaux de $\mathcal V$,
le signe de $\gamma$ change. L’évaluation du spectre sans ses étiquettes conserve
la paire non ordonnée, tandis que \eqref{iii:eq:barbero-trace} conserve la
contribution orientée. Si l’on emploie la trace normalisée
$\tau_2=\operatorname{Tr}/2$, la même formule s’écrit
$\gamma=-2\tau_2[\mathcal R\mathcal H(\psi(\mathcal V))]$.

En revanche, un changement simultané de représentation
$(\mathcal R,\mathcal V)\mapsto
(U\mathcal RU^{-1},U\mathcal VU^{-1})$ conserve la fonctionnelle :
le calcul spectral se transporte par conjugaison et la trace est
invariante. Échanger la réponse par rapport à une orientation fixe
et transporter conjointement la réponse et l’orientation sont donc
des opérations différentes.

Cette normalisation est compatible avec les amplifications de
\eqref{iii:eq:logvac}. Pour
$\mathcal V_m=\mathcal V\otimes I_{9^m}$ et
$\mathcal R_m=\mathcal R\otimes I_{9^m}$,
\begin{equation}
 \gamma=-\frac1{9^m}\operatorname{Tr}
 \left[\mathcal R_m\mathcal H\bigl(\psi(\mathcal V_m)\bigr)\right].
 \label{iii:eq:barbero-amplification}
\end{equation}
En effet, le calcul fonctionnel conserve le produit tensoriel et la trace
de $I_{9^m}$ est $9^m$. La multiplicité de la représentation se sépare
de l’évaluation orientée.

\subsection{Reconstruction à partir des coordonnées du vide}

Les identités \eqref{iii:eq:four-identities} permettent d’exprimer la même
fonctionnelle au moyen de l’impédance et de la vitesse normalisées :
\begin{align}
 \gamma={}&
 \mathcal H\!\left(\psi\!\left(
     \sqrt{\frac{\widehat Z}{\widehat c}}\right)\right)
 -\mathcal H\!\left(\psi\!\left(
     \frac1{\sqrt{\widehat Z\widehat c}}\right)\right).
 \label{iii:eq:barbero-zc}
\end{align}
Le domaine positif de cette reconstruction est décrit par
\begin{equation}
 1<\widehat Z\widehat c<\frac{16}{9},\qquad
 \frac9{16}<\frac{\widehat Z}{\widehat c}<1.
 \label{iii:eq:barbero-zc-domain}
\end{equation}
En effet, ces inégalités équivalent à
$3/4<r_\pm<1$ lorsqu’on substitue les racines positives de
\eqref{iii:eq:four-identities}. Pour l’orientation $y>0$, on conserve
en outre $r_-<r_+$, ou de manière équivalente $\widehat Z<1$. Les points du
domaine sont utilisés avec la compatibilité angulaire déjà établie par
le TPK ; l’inversion est une reconstruction ultérieure de ses canaux.
Les coordonnées sont les coordonnées internes de \eqref{iii:eq:four} : un changement
d’unités doit être défait au moyen de \eqref{iii:eq:undo-units} avant
d’appliquer $\psi$.

L’échange $y\mapsto-y$ agit comme
$(\widehat Z,\widehat c)\mapsto(\widehat Z^{-1},\widehat c)$ et
produit $\gamma\mapsto-\gamma$. Dans le cas équilibré $y=0$,
les deux canaux coïncident et \eqref{iii:eq:barbero-factorization} donne
$\gamma=0$. La vitesse normalisée conserve le produit des
canaux ; sa lecture conjointe avec l’impédance retrouve les deux
réponses et l’orientation nécessaire à l’évaluation de la fonctionnelle.
Dans ce cas équilibré, $\mathcal V=rI$ a un spectre dégénéré :
la marque $\mathcal R$ est conservée comme donnée de représentation, bien que
le scalaire $r$ ne distingue plus ses projecteurs. Pour un spectre simple,
en revanche,
$P_+=(\mathcal V-r_-I)/(r_+-r_-)$ et $P_-=I-P_+$.
